\documentclass[10pt]{beamer}

\usetheme{metropolis}
\definecolor{WiLabRed}{RGB}{197,18,48}
\setbeamercolor{frametitle}{fg=white,bg=WiLabRed}
\setbeamercolor{progress bar}{fg=WiLabRed!90}
\setbeamercolor{title separator}{fg=WiLabRed!90}
\setbeamercolor{progress bar in section page}{fg=WiLabRed!90}
\setbeamercolor{background canvas}{bg=white}
\setbeamercolor{alerted text}{fg=WiLabRed!90}

\usepackage{appendixnumberbeamer}

\usepackage{booktabs}
\usepackage[scale=2]{ccicons}

\usepackage{pgfplots}
\usepgfplotslibrary{dateplot}

\usepackage{xspace}
\newcommand{\themename}{\textbf{\textsc{metropolis}}\xspace}

\usepackage{marvosym}
%\usepackage{subfig}
\usepackage{graphicx}\graphicspath{{images/}}
\usepackage{subcaption}
\usepackage[framed]{./support/mcode}
\usepackage{listings}

\title{Lecture 03: Modeling Randomness in Communication Systems}
\subtitle{\textit{Software Defined Radio for Engineers} (Collins~\textit{et al.}), \textsection{3.3}}
\date{}
\author{\textbf{Alexander M. Wyglinski, Ph.D.}}
\institute{ \vspace*{1in}\hfill\includegraphics[height=1.125cm]{wilab_logo-A70916.eps} \qquad \includegraphics[height=1.125cm]{WPI_Inst_Prim_FulClr.eps}}
% \titlegraphic{\hfill\includegraphics[height=1.5cm]{logo.pdf}}

% Foot for all slides
\setbeamertemplate{frame footer}{\tiny \copyright~2018 by Alexander Wyglinski. This work is licensed under the Creative Commons Attribution-ShareAlike 4.0 International License. To view a copy of this license, visit http://creativecommons.org/licenses/by-sa/4.0/.}

\begin{document}

%\captionsetup[subfigure]{labelformat=empty}

%%%%%%%%%%%%%%%%%%%%%%%%%%%%%%%%%%%%%%%%%%%%%%%%%%%%%%%%%%

\maketitle




%%%%%%%%%%%%%%%%%%%%%%%%%%%%%%%%%%%%%%%%%%%%%%%%%%%%%%%%%%

\frame
{
  \frametitle{What Is Probability?}

  \begin{itemize}
   \item It is not about finding the exact value produced by a random event but rather accurately characterizing the random event itself and all its possible outcomes
   \item Two types of random events: \textit{discrete} and \textit{continuous}
   \item Probability theory involves \textit{mapping} physical phenomena to a mathematical framework such that we can employ a variety of useful tools for analysis
   \item We can represent a random event via the ``black box'' model
  \end{itemize}


}

%%%%%%%%%%%%%%%%%%%%%%%%%%%%%%%%%%%%%%%%%%%%%%%%%%%%%%%%%%

\frame
{
  \frametitle{Random Variables}

  \begin{itemize}
   \item Random variable ${X}$ is a function that maps an outcome $\omega$ from the sample space $\Omega$ to a real number $x$
   \item Sample space $\Omega$ is the domain of ${X}$ while $\Omega_x$ is the range
   \item How does this look like graphically?
  \end{itemize}

}


%%%%%%%%%%%%%%%%%%%%%%%%%%%%%%%%%%%%%%%%%%%%%%%%%%%%%%%%%%

\frame
{
  \frametitle{Probability}

  \begin{itemize}
   \item Probability that discrete random variable ${X}$ produces an output within $[a,b)$ is given by:
   \begin{equation}
    P(a\le{X}<b)=\sum\limits_{i=a}^{b-1}p_X(i)
   \end{equation}
   where $p_X(i)$ is the probability mass function (PMF)
  \item Probability that continuous random variable ${X}$ produces an output within $[a,b)$ is given by:
   \begin{equation}
    P(a\le{X}<b)=\int\limits_{a}^{b}f(t)dt
   \end{equation}
   where $f(t)$ is the probability density function (PDF)
  \end{itemize}


  }


%%%%%%%%%%%%%%%%%%%%%%%%%%%%%%%%%%%%%%%%%%%%%%%%%%%%%%%%%%

\frame
{
  \frametitle{Expectation}

  \begin{itemize}
   \item Probability is about characterizing random events and their behavior
   \item One such mathematical tool is \textit{Expectation}
   \begin{itemize}
    \item It can be roughly described as a weighted average of all possible outcomes of a random event
   \end{itemize}
   \item We define Expectation as:
   \begin{equation}
    E[g(X)]=\int\limits_{-\infty}^{+\infty}g(x)f(x)dx
   \end{equation}
    \item We use Expectation to give us some understanding about the behavior of a random variable
  \end{itemize}


  }



%%%%%%%%%%%%%%%%%%%%%%%%%%%%%%%%%%%%%%%%%%%%%%%%%%%%%%%%%%

\frame
{
  \frametitle{CDFs and PDFs}

  \begin{itemize}
   \item To \textbf{perfectly} describe a random event, we use \textit{Cumulative Distribution Functions} (CDFs) and \textit{Probability Density Functions} (PDFs)
   \item Mathematically, this is described as:
   \begin{equation}
    F_X(x)=P(X\le{x})=\int\limits_{-\infty}^{x}f(t)dt
   \end{equation}
   \item Here is a useful trick for calculating the probability $P(a\le{X}<b)$:
   \begin{equation}
    P(a\le{X}<b)=P(X<b)-P(X<a)=F_X(b)-F_X(a)
   \end{equation}
  \end{itemize}

  }


%%%%%%%%%%%%%%%%%%%%%%%%%%%%%%%%%%%%%%%%%%%%%%%%%%%%%%%%%%

\frame
{
  \frametitle{Important Communications-Related PDFs}

  \begin{itemize}
   \item Uniform:
   \begin{equation}
    f(x)=\begin{cases}
          \frac{1}{(b-a)},&a\le{x}<{b}\\
          0,&\mathrm{otherwise}
         \end{cases}
   \end{equation}
   \item Gaussian:
   \begin{equation}
    f(x)=\frac{1}{\sqrt{2\pi{\sigma^2}}}e^{-((x-\mu)/\sigma)^2/2}
   \end{equation}
  \end{itemize}

  }



%%%%%%%%%%%%%%%%%%%%%%%%%%%%%%%%%%%%%%%%%%%%%%%%%%%%%%%%%%

\frame
{
  \frametitle{Q-Function}

  \begin{itemize}
   \item \textit{Q function} is a convenient way to express right-tail probabilities for Gaussian random variables
   \item Mathematically, this can be expressed as:
   \begin{eqnarray}
    Q(x)&=1-F_X(x)=1-P(X\le{x})\\
    &=P(X>x)=\frac{1}{\sqrt{2\pi}}\int\limits_{x}^{\infty}e^{-t^2/2}dt
   \end{eqnarray}
  \end{itemize}

  }


%%%%%%%%%%%%%%%%%%%%%%%%%%%%%%%%%%%%%%%%%%%%%%%%%%%%%%%%%%

\frame
{
  \frametitle{Where Does Randomness Occur?}

  \begin{itemize}
   \item Several locations where randomness occurs in the communications system includes:
   \begin{itemize}
    \item Transmission channel
    \item Data generation at information source
    \item Clock jitter
    \item Processing latency
    \item $\ldots$
   \end{itemize}

  \end{itemize}


  }


\end{document}
